\lyseis{Α2.4}{Πολλαπλασιασμός ακεραίων}
\sol{1} α) $6$ \quad β) $-24$ \quad γ) $42$ \quad δ) $-25$ \quad ε) $0$ \quad στ) $-207$.
\sol{2} α) θετικός. β) αρνητικός. γ) αρνητικών. δ) ετερόσημοι.
\sol{3} α) Πέντε αρνητικοί (περιττό πλήθος): $-1$. β) Τέσσερις αρνητικοί ($-1$, $-2$, $-4$, $-5$), άρτιο πλήθος, και $1\cdot2\cdot3\cdot4\cdot5=120$: $+120$. γ) Τρεις αρνητικοί, και $5\cdot10\cdot5\cdot10\cdot2=5000$: $-5000$.
\sol{4} Γραμμή $-5$: $-15$, $20$, $30$, $0$. Γραμμή $-1$: $-3$, $4$, $6$, $0$. Γραμμή $+2$: $6$, $-8$, $-12$, $0$. Γραμμή $7$: $21$, $-28$, $-42$, $0$.
\sol{5} α) $(+3)\cdot(-4)=-12$ και $-12=(-2)\cdot6$, άρα $6$. β) $(-5)\cdot(-6)=30=3\cdot10$, άρα $10$. γ) $9\cdot8=72=(-6)\cdot(-12)$, άρα $-12$. δ) $(-3)\cdot8=-24=12\cdot(-2)$, άρα $-2$.
\sol{6} α) $2\cdot(-2)+4\cdot5=-4+20=16$. β) $(-9)\cdot(-1)=9$ (γιατί $-7+11-13=-9$ και $-5+4=-1$). γ) $-18+(-4)-40=-62$ (ο όρος $2\cdot(-5)\cdot(-4)=40$ αφαιρείται). δ) $3\cdot[-8-(-10)]=3\cdot2=6$.
\sol{7} α) Σωστό. β) Λάθος: $+12$. γ) Σωστό. δ) Λάθος: οι $100$ αρνητικοί είναι άρτιο πλήθος, άρα $+1$. ε) Σωστό (αντιμεταθετική). στ) Λάθος: $(-2)\cdot(-3)=6$. ζ) Σωστό.
\sol{8} $\alpha\beta=-4$, $\beta\gamma=-20$, $\alpha\beta\gamma=(-4)\cdot(-5)=20$, $\alpha\gamma-\beta=5-4=1$, $\alpha-\beta\gamma=-1-(-20)=19$.
\sol{9} α) $4\cdot(-3)=-12$, δηλαδή $-12$°. β) $(-3)\cdot(-2)=+6$, δηλαδή $+6$° (είχε $6$ βαθμούς περισσότερο).
\sol{10} Ηλίας: $7\cdot5+4\cdot(-2)=35-8=27$ πόντους. Στέλιος: $3\cdot5+9\cdot(-2)=15-18=-3$ πόντους.
\sol{11} $6+5\cdot(-2)=6-10=-4$. Βρίσκεται στον όροφο $-4$ (τέταρτο υπόγειο).
\sol{12} Μετά από $6$ μέρες: $6\cdot(-15)=-90$, δηλαδή $90$ λίτρα λιγότερα. Πριν από $4$ μέρες: $(-15)\cdot(-4)=+60$, δηλαδή υπήρχαν $60$ λίτρα περισσότερα.
\sol{13} $(-6)\cdot(-1)=6$ και $(-6)\cdot(-2)=12$. Κάθε γραμμή είναι κατά $6$ μεγαλύτερη από την προηγούμενη, άρα το γινόμενο δύο αρνητικών είναι θετικό.
\sol{14} Αν $\alpha\beta>0$, οι $\alpha,\,\beta$ είναι ομόσημοι (και οι δύο θετικοί ή και οι δύο αρνητικοί). Αν $\alpha\beta<0$, είναι ετερόσημοι.
\sol{15} Άρτιο πλήθος αρνητικών: $0$, $2$ ή $4$.
\sol{16} Ο κανόνας «μείον επί μείον κάνει συν» ισχύει μόνο στον πολλαπλασιασμό. Στην πρόσθεση δύο αρνητικοί δίνουν αρνητικό: $(-3)+(-4)=-7$, ενώ $(-3)\cdot(-4)=+12$.
\sol{17} Οι αριθμοί είναι ετερόσημοι. Δοκιμάζω: $4\cdot(-3)=-12$ και $4+(-3)=1$. Αυτό είναι το ζεύγος $4$ και $-3$.
\sol{18} Ελεύθερη κατασκευή. Η γραμμή και η στήλη του $0$ έχουν παντού $0$. Στα κελιά με ομόσημους αριθμούς το γινόμενο είναι θετικό και στα κελιά με ετερόσημους αρνητικό.
\sol{19} Ελεύθερη απάντηση. Π.χ. αν μια ταινία που δείχνει κάποιον να κατεβαίνει $2$ ορόφους την ώρα γυρίσει ανάποδα (αρνητικός χρόνος), δείχνει να ανεβαίνει: $(-2)\cdot(-1)=+2$.
